\subsection{A sharper bound for component-sharing types}
\label{sec:types}

The recovered grading also improves the finite-type estimate used by the
earlier structural-compression report. This is an improvement to a
sufficient parameter regime, beyond the faster composition kernel.

For integers $m\geq0$ and $b\geq1$, let
\[
 B_m=\binom{m}{\floor{m/2}},\qquad M(2m)=(2m-1)!!,
\]
with $M(0)=1$. A homogeneous element of a $k$-circle coefficient space
belongs to one of its $k+1$ dot-degree layers. Each layer has dimension
at most $B_m$ for $k\leq m$. Hence the number of possible coefficients,
including zero, is at most $H_m=(m+1)2^{B_m}$. This union bound
intentionally overcounts zero.

\begin{theorem}[Homogeneous component-type bound]
\label{thm:types}
At a fixed frontier of $2m$ labeled points, the number of serialized
connected differential components with at most $b$ objects, homogeneous
entries, and homological degrees normalized by their minimum is at most
\begin{equation}
 K(m,b)=\sum_{s=1}^{b}s^s M(2m)^s H_m^{s^2}.
 \label{eq:types}
\end{equation}
Consequently
\begin{equation}
 \log K(m,b)=O\bigl(b^2 B_m+b m\log(m+1)+b\log(b+1)\bigr)
 =O(b^2 B_m).
 \label{eq:typecost}
\end{equation}
\end{theorem}
\begin{proof}
For a connected component of $s$ objects, every nonzero differential
changes homological degree by one. An undirected path of at most $s-1$
edges joins a minimum-degree vertex to any other vertex. The normalized
degrees therefore lie in $\{0,\ldots,s-1\}$, giving at most $s^s$
assignments. There are at most $M(2m)^s$ choices of matching labels. For
each ordered pair of vertices, the homogeneous coefficient belongs to a
set of at most $H_m$ possibilities. This gives $H_m^{s^2}$ choices.
We have ignored the required degree difference, connectedness, and
$d^2=0$, so the count is an upper bound. It already counts ordered
serializations; no graph-isomorphism algorithm is assumed.

Take logarithms and bound the sum by $b$ times its largest crude term.
The estimates $\log M(2m)=O(m\log(m+1))$ and the definition of $H_m$
give the first expression. Finally, $m\log(m+1)=O(B_m)$,
$\log(m+1)=O(B_m)$, and $\log(b+1)=O(b)$ uniformly over nonnegative
integers, with the finitely many small values absorbed into constants.
This gives the second expression.
\end{proof}

The keys in the earlier component-sharing code do not store quantum
shifts. This does not invalidate the count: Proposition~\ref{prop:grading}
ensures that each observed coefficient is homogeneous. Forgetting a
hidden grading cannot increase the number of structures compared with
the set of possible graded structures. There is also a potentially
tighter count that assigns normalized quantum shifts to vertices.
Across a nonzero edge their difference is
$m-k+2d\in[0,2m]$. In a connected $s$-vertex component their range is
therefore at most $2m(s-1)$. Once both matching and quantum shifts are
fixed, every ordered pair has a uniquely determined dot-degree layer.
This yields the alternative upper bound
\[
 \sum_{s=1}^{b}s^s[2m(s-1)+1]^s M(2m)^s 2^{s^2 B_m}.
\]
Its leading exponent has the same order as~\eqref{eq:typecost}.

\begin{corollary}[Improved sufficient regime for exact sharing]
\label{cor:sharing}
Consider the exact component-sharing scanner from report 07, using the
present coefficient kernel. Suppose every reduced connected component
encountered in an execution has at most $b$ objects and the frontier has
at most $2m$ points. Then a sufficient condition for
$n^{O(\log n)}$ time is
\begin{equation}
 b^2\binom{m}{\floor{m/2}}=O((\log(n+2))^2).
 \label{eq:sharingregime}
\end{equation}
\end{corollary}
\begin{proof}
There are at most $K(m,b)$ stored component representatives, each with
at most $b$ objects. Tensoring a representative with one crossing creates
only a constant factor more objects before cancellation. Distinct parent
summands stay independent under this operation, so the representation
has polynomial size in $bK(m,b)$. Its homological-shift weights have
degree $O(n)$ and coefficients of $O(n)$ bits: the full cube has at most
exponentially many generators in $n$. Saturated decision weights are
smaller still. The resource envelope is therefore polynomial in
$n,b,K(m,b)$ times $\poly(m)2^m$. Insert
Theorem~\ref{thm:types} and~\eqref{eq:sharingregime}.
\end{proof}

Since $B_m=\Theta(2^m/\sqrt{m+1})$, the sufficient leading term improves
from $b^2 2^m$ in the generic ungraded type count to
$b^2 2^m/\sqrt{m+1}$. For bounded $b$, put
$L=\log_2(n+2)$ and $\ell=\log_2 L$. The corresponding frontier scale
is enlarged from $w\leq4\ell+O(1)$ to
\[
 w\leq4\ell+\log_2\ell+O(1).
\]
The additive constant can depend on the chosen bound for $b$.
This is a modest asymptotic refinement of the earlier regime. The
unresolved task is still to prove useful bounds on actual connected
component sizes for broad knot families. The archive does not claim a
new benchmark of the report-07 scanner with every recent branch merged.
